\section{The operational clauses and their source locators}\label{app:clauses}
We use the printed pagination of Marletto's May 2017 version, arXiv:1608.02625v5~\cite{marletto}. A locator identifies the source requirement. The following pair-task formulation specifies where we retain a requirement and where we restrict or replace it. Its realization is proved in Appendices~\ref{app:permutation}--\ref{app:media}.

\subsection{Attributes, dynamics and the constructor}
A thermodynamic attribute here is a specified density operator on a finite substrate. Independent attributes compose by tensor product; correlated states are also attributes of the whole composite. This specializes the source's set-valued attributes to its density-operator and thermal-state examples in Sec.~2, pp.~11--12. Finite rational Hamiltonians, the admissible state domain and the gate class are subsidiary specifications.

The law acting on a state is deterministic. In the completion it is an energy-preserving unitary, with no random choice of gate. This is our subsidiary dynamics within the nonprobabilistic task framework of principle I, p.~14; determinism alone is not that principle. Locality, principle II, p.~15, is implemented by requiring that an operation confined to one factor does not change the other factor's Heisenberg descriptors. It does not assert that an entangled state is a pair of independent marginal states.

The constructor contract is~\eqref{eq:closed}: the whole memory marginal returns exactly at each visit, and the system output approaches the target. It is an explicit interpretation of retained ability in Sec.~2, pp.~13--14, and Sec.~2.1, pp.~16--18, not a return condition stated there. Returning internal memory registers separately is not sufficient. A fresh input is independent of the memory and history; the required output is its marginal, not a fresh independent sample.

We restrict the composition requirement in principle III, p.~15, to serial and parallel composition of thermodynamic pair tasks. The completion also supplies actual multiple-input information and work gates. For each such gate, a single circuit must handle every promised input. Separate existence of its branches is not a substitute. The restricted clause does not identify a composite constructor with a serial connection of arbitrary catalysts whose memories may become correlated.

Stable storage and preparation from a sufficiently large finite supply of generic resources implement the requirements used from principles IV and V in Sec.~2.1, pp.~17--19. Resources consumed in preparation are recorded as side effects. They are not supplied for free during a closed task. Information variables have perfectly distinguishable labels that admit copying and permutations with counted side effects. Their products are information variables, and a finite pairwise-distinguishable family is jointly distinguishable. These are the information requirements used from Sec.~3, principles VI and VII, pp.~20--21.

Energy adds on independent factors and is conserved in every microscopic gate, as required by VIII, p.~25. Its spectrum is bounded on each individual substrate, as on pp.~25--26. A sequence of larger constructors is not one substrate with unbounded energy. We introduce no additional conserved charge.

\subsection{Work and the calibration closure}
A work variable here is a finite, equally spaced list of known pure energy configurations, with a nonzero rational spacing. One may embed the list in a larger substrate. On a composite we use one known product configuration for each distinct energy in the list. Work sharpness supplements the mechanical-coordinate interpretation in Sec.~5, pp.~27--29 and 33: a change of internal randomness is not accepted as a work displacement.

Conditions 1, 2a and 2b in Sec.~5, pp.~27--28, are retained for these lists. Distinct isolated work levels cannot be interconverted without a side effect. Adjacent forward steps are is-like: combining either with the transpose of the other is possible without an external side effect. Every adjacent pair admits the source's two-branch swap, using a replica initialized in one member of that pair. The work attributes used as side effects also admit swaps with either replica initialization, as required for the source's subset $W'$ on p.~28. We choose guarded internal levels for this purpose. Guard levels belong to a finite register, not an infinite reservoir.

We replace condition 2c, p.~28, by a weaker rule at the stated physical factorization. Each changing factor of a factorized mechanical pair is a known pure-state mechanical transition; unchanged factors are identities. We do not require every component to be a work variable. In the source, two-level systems are instead quasi-work media, not work media (p.~33); the replacement is not a correction of that distinction. We also stipulate that appending a fixed pure spectator preserves a work list. Its dimension and energy span are included when it is absorbed into the work substrate, as in Theorem~\ref{thm:calibration}.

Work interoperability I, p.~32, equation (4), requires the same two-branch swap between matched work lists on different substrates, with the matched initial level in the ancillary list. The printed formula's cross-substrate initial-state label is read through this matching. For part (ii), we choose one pure representative of each distinct composite energy rather than identify physically different configurations as one pure state. This is a refinement of the source's composite-variable convention. For equal-spacing matched lists these representatives form an equally spaced list. Configurations omitted only because their energies are repeated remain available and can be interchanged reversibly.

Adiabatic possibility uses work side effects only, as in Sec.~5, pp.~33--34. We replace a fixed exactly calibrated side-effect pair by~\eqref{eq:adiabatic}, with dimension and energy span bounded uniformly as accuracy improves. This additional accuracy convention permits rational work gaps to converge to a mean-energy difference without letting small probability tails carry finite unreported entropy or energy. It does not demand a common bound across different pair tasks, nor a bound on constructor memory.

\subsection{Heat, the working interval and comparison}
A heat variable is a designated family on a nonconstant finite-energy substrate. Distinct members have exactly one adiabatically possible direction, retaining the quasi-heat requirement on p.~34. The four tests on p.~35 are used with the following restrictions and replacement.

First, no member can be distinguished from any other supplied thermodynamic attribute with arbitrarily small error by one-use processing with an independently prepared apparatus. Second, its identical-copy family is again quasi-heat. Third, any two members admit reversible adiabatic equalization to two copies of a member of the same family. We use the bounded calibration closure for this work-assisted equalization. A single fixed rational work witness would fail for some finite Gibbs pairs; Appendix~\ref{sec:rationalwitness} proves this obstruction. Fourth, no heat attribute, when supplied as the initial ancillary state, can enable the upward branch of a sharp work swap. The ancillary output may be arbitrary. This state-specific prohibition replaces the literal substrate-wide prohibition on being a work-like ancilla. The heat-state discussion on pp.~36--37 motivates it, but does not make the two formulations identical.

We restrict heat interoperability IX, p.~38, to whole designated families. A correspondence must make every pair task is-like to its corresponding pair task. The matched diagonal product family must then satisfy all four heat tests. Matching a selected finite pair is not taken to establish interoperability of every temperature in two continua. Thus the clause used here is narrower than a requirement on all source heat variables.

We weaken existence principle II, p.~39, in two respects. Heat counterparts are required only for work energies in the declared thermal working interval, and work attributes with the same energy may share a single counterpart. The source explicitly includes equal-energy work pairs and then requires the corresponding heat task to have one possible direction and an impossible transpose. A shared counterpart gives an identity task and does not satisfy that strict requirement. It is therefore a replacement of the source clause, not a clarification of the word ``pair.''

The zeroth law X, p.~42, is retained on this domain: any thermodynamic state can reach any designated heat attribute at the same energy without a work side effect. The interval avoids demanding a full-rank heat state at an extremal energy of a nonconstant finite Hamiltonian. It is part of the modified formulation, not a derived universal temperature range.

\begin{lemma}\label{lem:injectivity}
Zeroth law X and strict accessibility imply that energy is injective on each heat variable.
\end{lemma}
\begin{proof}
If $h$ and $k$ in that variable have equal energy, X gives both $h\rightsquigarrow k$ and $k\rightsquigarrow h$. A closed task is adiabatic with an unchanged work register. The strict requirement for distinct heat attributes therefore gives $h=k$.
\end{proof}
Within this pair-task formulation, retaining X and strict accessibility rules out a distinct same-energy heat pair. The common-counterpart convention changes the existence clause to respect this consequence.

We impose comparison for every pair on the same finite composite. The source obtains it under a simplifying single-law assumption on p.~40, while footnote~12 says that it is not required in that treatment. Here it is an explicit assumption. Kelvin's prohibition excludes a closed heat-to-work conversion (Sec.~7, pp.~42--44). The completion additionally prohibits positive sharp work extraction from a single positive-temperature Gibbs state with arbitrary final state.

Marletto asserts on p.~40 that adiabatic is-like is an equivalence relation, without proving transitivity there. We do not include that assertion as an axiom. Appendix~\ref{app:cancellation} identifies it with cancellation under comparison. The real additive entropy label introduced on pp.~40--41 is likewise not assumed. At fixed energy the completion has one oriented obstruction, namely entropy decrease; comparison is a separate assumption in the general theorem, not a restatement of the source's single-law-of-impotence simplification.
